\subsection{Relative invariant of two lattices}

We keep the notation and hypotheses of nos. 2 to 5. Let V be a vector space of finite rank n over K and M a lattice of V with respect to A. Let W be the exterior power \(\bigwedge^{n}V\) , which is a vector space of rank 1 over K and let \(M_{W}\) denote the lattice of W generated by the image of \(M^{n}\) under the canonical mapping \(V^{n}\to\bigwedge^{n}V\) (no. 1, Proposition 3 (iii); note that \(M_{W}\) is not necessarily isomorphic to \(\bigwedge^{n}M\) (Algebra, Chapter III, § 5, Exercise 9)). If e is a basis of W over K, we may write \(M_{W}=a.e\) , where a is a fractional ideal \(\neq0\) of A.

Let \(\mathbf{M}'\) be another lattice of \(\mathbf{V}\) and let us write \(\mathbf{M}_{\mathbf{W}}' = \mathfrak{a}'\) .e, where \(\mathfrak{a}'\) is a fractional ideal \(\neq 0\) of \(\mathbf{A}\) ; the divisor \(\operatorname{div}(\mathfrak{a}) - \operatorname{div}(\mathfrak{a}')\) does not depend on the choice of basis \(e\) of \(\mathbf{W}\) , \(\mathfrak{a}\) and \(\mathfrak{a}'\) being multiplied by the same element of \(\mathbf{K}^*\) when the basis is changed; we shall write \(\chi(\mathbf{M}, \mathbf{M}') = \operatorname{div}(\mathfrak{a}) - \operatorname{div}(\mathfrak{a}')\) and say that this divisor is the relative invariant of \(\mathbf{M}'\) with respect to \(\mathbf{M}\) . Clearly, if \(\mathbf{M}, \mathbf{M}', \mathbf{M}''\) are three lattices of \(\mathbf{V}\) , then:

\[
\chi (\mathbf {M}, \mathbf {M} ^ {\prime}) + \chi (\mathbf {M} ^ {\prime}, \mathbf {M} ^ {\prime \prime}) + \chi (\mathbf {M} ^ {\prime \prime}, \mathbf {M}) = 0\tag{3}
\]

\[
\chi (\mathbf {M}, \mathbf {M} ^ {\prime}) + \chi (\mathbf {M} ^ {\prime}, \mathbf {M}) = 0\tag{4}.
\]

For all \(p \in P\) , it follows immediately from the definitions that \((M_{W})_p = (M_p)_{W}\) ; moreover, since \(M_p\) is then a free \(A_p\) -module since \(A_p\) is a principal ideal domain, a basis of \(M_p\) over \(A_p\) is a basis of \(V\) over \(K\) , hence \((M_p)_W = \bigwedge^n (M_p)\) (Chapter II, § 2, no. 8) and the fractional ideal \(a_p = aA_p\) is a principal ideal domain. If we set \(a_p = p^{n_p} A_p\) , \(a' = p^{n'_{p}} A_p\) , then:

\[
\chi (\mathbf {M}, \mathbf {M} ^ {\prime}) = \sum_ {p \in P} (n _ {p} - n _ {p} ^ {\prime}). p,
\]

which may also be written as:

\[
\chi (\mathbf {M}, \mathbf {M} ^ {\prime}) = \sum_ {p \in P} \chi (\mathbf {M} _ {p}, \mathbf {M} _ {p} ^ {\prime})\tag{5}
\]

identifying \(\mathbf{D}(\mathbf{A}_{\mathfrak{p}})\) with the sub-Z-module of \(\mathbf{D}(\mathbf{A})\) generated by p.

PROPOSITION 13. Let M be a lattice of V and u a K-automorphism of V. Then:

\[
- \chi (\mathbf {M}, u (\mathbf {M})) = \operatorname{div} (\det (u)).\tag{6}
\]

For all \(p \in P\) , then \(\bigwedge^{n} (u(M)_{p}) = \bigwedge^{n} (u(M_{p}))\) ; if \((e_{i})_{1 \leqslant i \leqslant n}\) is a basis of \(M_{p}\) , then

\[
\bigwedge^ {n} \left(\mathrm{M} _ {p}\right) = \mathrm{A} _ {p}. e _ {1} \wedge e _ {2} \wedge \dots \wedge e _ {n},
\]

and \(\bigwedge^{n}(u(\mathbf{M}_{\mathfrak{p}})) = \mathbf{A}_{\mathfrak{p}}.\det (u)e_{1}\wedge e_{2}\wedge \dots \wedge e_{n}\) , whence the proposition by virtue of formula (5).

PROPOSITION 14. If M, M' are two lattices of V such that M' ⊂ M, M/M' is a finitely generated torsion A-module and:

\[
\chi (\mathbf {M}, \mathbf {M} ^ {\prime}) = - \chi (\mathbf {M} / \mathbf {M} ^ {\prime}).\tag{7}
\]

Clearly \(\mathbf{M} / \mathbf{M}' \subset \mathbf{V} / \mathbf{M}'\) is a finitely generated torsion module; on the other hand, for all \(\mathfrak{p} \in \mathbb{P}\) , we know (Algebra, Chapter VII, § 4, no. 2, Theorem 1) that there exist bases \((e_i)_{1 \leqslant i \leqslant n}\) of \(\mathbf{M}_{\mathfrak{p}}\) and \((e_i')_{1 \leqslant i \leqslant n}\) of \(\mathbf{M}_{\mathfrak{p}}'\) such that \(e_i' = \pi^{\nu_i} e_i\) for \(1 \leqslant i \leqslant n\) and integers \(\nu_i \geqslant 0\) , \(\pi\) being a uniformizer of \(\mathbf{A}_{\mathfrak{p}}\) . Therefore (in the notation introduced above) \(n_{\mathfrak{p}}' - n_{\mathfrak{p}} = \sum_{i=1}^{n} \nu_i\) ; and on the other hand, \((\mathbf{M} / \mathbf{M}')_{\mathfrak{p}} = \mathbf{M}_{\mathfrak{p}} / \mathbf{M}_{\mathfrak{p}}'\) is isomorphic to the torsion \(\mathbf{A}_{\mathfrak{p}}\) -module \(\bigoplus_{i=1}^{n} \mathbf{A}_{\mathfrak{p}} / p^{\nu_i} \mathbf{A}_{\mathfrak{p}}\) and hence its length is \(\sum_{i=1}^{n} \nu_i\) , which proves the proposition, in view of (5) and Definition 4 of no. 5.

COROLLARY. Let \(L_{1}\) , \(L_{2}\) be two free A-modules of the same rank n and let \(f: L_{1} \to L_{2}\) be a homomorphism. Let U be the matrix of f with respect to bases of \(L_{1}\) and \(L_{2}\) . For Coker(f) to be a torsion A-module, it is necessary and sufficient that \(\det(U) \neq 0\) and then:

\[
\chi (\operatorname{Coker} (f)) = \operatorname{div} (\det (U)).\tag{8}
\]

\(L_{1}\) and \(L_{2}\) can be considered as lattices in \(V_{1} = L_{1} \otimes_{A} K\) and \(V_{2} = L_{2} \otimes_{A} K\) respectively, f extending to a K-homomorphism \(f_{(K)}: V_{1} \to V_{2}\) . Then

\[
\left(\operatorname{Coker} (f)\right) _ {(\mathbf {K})} = \operatorname{Coker} \left(f _ {(\mathbf {K})}\right)
\]

and to say that \(\operatorname{Coker}(f)\) is a torsion A-module means that \(\operatorname{Coker}(f_{(\mathbf{K})}) = 0\) ; now, it amounts to the same to say that \(f_{(\mathbf{K})}\) is surjective or that \(\det(U) \neq 0\) , whence the first assertion. On the other hand, we may write \(f(\mathbf{L}_1) = u(\mathbf{L}_2)\) , where \(u\) is an endomorphism of \(\mathbf{L}_2\) of determinant \(\det(U)\) ; as

\[
\operatorname{Coker} (f) = \mathrm{L} _ {2} / u (\mathrm{L} _ {2}),
\]

formula (8) follows from (7) and (6).

Example. If A = Z, the divisor group of A is identified with the multiplicative group \(Q_{+}^{*}\) of rational numbers \textgreater0. For every finite commutative group T, \(\chi(T)\) is the order of T; the above corollary shows that the order of the group \(\operatorname{Coker}(f)\) is equal to the absolute value of \(\det(U)\) (cf.~Algebra, Chapter VII, §4, no. 7, Corollary 3 to Theorem 3).
% semantic-fragments: legacy-input-seam
\relax{}
